\section{Introduction} \label{sec:51-research03-introduction}

\begin{foldable} % Text DD
  The systemic costs that motivate dataset compression sharpen at the corpus scale: at extreme compression, the residual budget has to be spent on the samples that contribute the most to learning.
  Yet, adapting image \ac{DD} techniques~\cite{wang2018dataset,zhao2021dataset,zhao2023dataset} to text faces two concrete obstacles: non-differentiable token decoding blocks gradient-based synthesis, and embedding spaces are tightly coupled to specific architectures.
\end{foldable}

\begin{foldable} % Baseline text DD methods --> gap --> desiderata
  Existing text \ac{DD} methods~\cite{sucholutsky2021soft,li2021data,maekawa2023dataset,maekawa2024dilm,tao2024textual} fall short in at least one of the following: at extreme compression, uniform weighting exhausts the budget on uninformative samples; their objectives lack task alignment or global optimization; and their outputs are embeddings, neither auditable nor transferable.
  We survey this landscape in \S\ref{sec:26-background-datasetdistill-textdd} and show that no prior method simultaneously addresses all of the above.
  These gaps motivate three desiderata for TAKE:{
    (1)~\textit{knowledge-based weighting} --- weighting samples by their task-aligned downstream contribution, directing the budget toward the most informative ones;
    (2)~a \textit{dataset-level distillation objective} --- optimizing over the full training distribution rather than batch-by-batch, ensuring global coverage;
    and (3)~\textit{human-readable output} --- a distilled corpus that is directly inspectable, auditable, and transferable across architectures.
  }
\end{foldable}

\begin{table}[ht]
  \centering
  \begin{threeparttable}
    \caption{Landscape of text dataset distillation methods.}
    \label{tab:landscape-textdd-methods}
    \setlength{\tabcolsep}{8pt}
    \begin{tabular}{l c c c c c c}
      \toprule
      \textbf{Gap}         & SLDD         & DDTC         & DDAL         & DiLM         & DaLLME       & \textbf{TAKE} \\
      \midrule
      Task alignment       & $\checkmark$ & $\checkmark$ & $\triangle$  & $\triangle$  & $-$          & $\checkmark$ \\
      Global optimization  & $\triangle$  & $\triangle$  & $\triangle$  & $\triangle$  & $\triangle$  & $\checkmark$ \\
      Interpretability     & $\triangle$  & $-$          & $-$          & $\checkmark$ & $\checkmark$ & $\checkmark$ \\
      Importance weighting & $-$          & $-$          & $-$          & $-$          & $-$          & $\checkmark$ \\
      \bottomrule
      \multicolumn{7}{l}{$-$: absent, $\triangle$: partial, $\checkmark$: addressed}
    \end{tabular}
  \end{threeparttable}
\end{table}

\begin{foldable} % HSB framing and positioning
  The \textit{importance weighting} row in Table~\ref{tab:landscape-textdd-methods} stands out: every prior method leaves it open.
  The natural fix is influence functions~\cite{koh2017understanding}, which quantify each sample's counterfactual effect on the downstream loss.
  Scoring at a single checkpoint, however, introduces the \textit{\acl{HSB}} (\ac{HSB}): at convergence, clean and moderate samples have near-zero gradient by definition, so noisy or hard samples dominate the influence scores, exactly inverting the ranking a robust distilled corpus needs.
  In standard supervised training the bias can often be absorbed by regularization, but in the extreme low-data regime of \ac{DD} it amplifies instead: with only a handful of examples per class to allocate, biased selection actively suppresses the informative samples the surrogate depends on.
  To our knowledge, no prior text \ac{DD} method has identified or corrected for the \ac{HSB}, which is why the table's last row remains uniformly empty.
\end{foldable}

\begin{foldable} % Our method TAKE and contributions
  We propose \textbf{Trajectory-Aware Knowledge Estimation} (\textbf{TAKE}), the first text \ac{DD} method to satisfy all three desiderata.
  TAKE corrects the \ac{HSB} by integrating influence scores along the full training trajectory into a per-sample knowledge score that captures clean and moderate samples while down-weighting noisy ones (\S\ref{sec:52-research03-theory}).
  Then, TAKE fine-tunes a \ac{LLM} to generate a pool of human-readable candidate instances.
  The knowledge scores and synthetic candidates jointly feed a discrete \ac{OT} objective that globally selects the distilled corpus.
  Although \ac{OT} has been widely applied in \ac{NLP} and machine learning, to our knowledge it has not been applied to text \ac{DD}, nor combined with a knowledge-reweighted source distribution in any \ac{DD} setting; TAKE fills both gaps.
  Concretely, we contribute:
  \begin{itemize}
    \item {
        \textbf{Theory:} We formalize the \ac{HSB} in the text \ac{DD} setting, prove that single-checkpoint influence scores are biased toward noisy samples, and derive a formal gap bound for knowledge-reweighted distribution matching.
      }
    \item {
        \textbf{Method:} TAKE is the first text \ac{DD} method to jointly address sample weighting and distributional coverage, via trajectory-integrated knowledge scores and a discrete \ac{OT} objective aligned with the \ac{DD} goal.
      }
    \item {
        \textbf{Empirical:} TAKE matches or exceeds prior text \ac{DD} methods across six language benchmarks at extreme compression (0.1\% or 20 instances/class), producing human-readable distilled corpora that generalize across diverse backbone families.
      }
  \end{itemize}
\end{foldable}

\begin{foldable} % Roadmap
  The chapter follows the order in which the three desiderata are settled.
  \S\ref{sec:52-research03-theory} casts text \ac{DD} as reweighted distribution matching problem and shows where the naive instantiation fails.
  \S\ref{sec:53-research03-method} then builds the TAKE pipeline: gradient-based trajectory influence and its knowledge score, synthetic candidate generation, and distributional prototype selection via discrete \ac{OT}.
  \S\ref{sec:54-research03-experiments} evaluates TAKE on classification and natural language inference benchmarks at extreme compression and confirms that each component is load-bearing.
  \S\ref{sec:55-research03-limitations} discusses limitations and ethical considerations, and \S\ref{sec:56-research03-conclusion} closes the chapter.
\end{foldable}
